\documentclass[10pt,a4paper]{amsart}
\usepackage{amsmath,amssymb,amsthm,graphicx}
% Retained preamble for the Stacks Project chapter divisors.tex (Divisors).
% The chapter file pinned-original-source.tex carries only \input{preamble} on its first line, so no
% preamble text lives inside the chapter itself.  The Stacks Project uses one single shared preamble for
% every chapter; the snapshot stages it verbatim in sources/stacks-chow/preamble.tex.  The block below is
% copied mechanically, line for line, from that file's lines 14-18 (the Xy-pic setup) and lines 154-217
% (the native theorem declarations together with the mathematical macros).  Nothing was added, removed or
% redefined.  The wrapper class and the amsmath/amssymb/amsthm/graphicx base packages are supplied by the
% assembler, so the class/\externaldocument part of the shared preamble is not repeated here.
% For commutative diagrams we use Xy-pic
\usepackage[all]{xy}

% We use 2cell for 2-commutative diagrams.
\xyoption{2cell}
% Theorem environments.
%
\theoremstyle{plain}
\newtheorem{theorem}[subsection]{Theorem}
\newtheorem{proposition}[subsection]{Proposition}
\newtheorem{lemma}[subsection]{Lemma}

\theoremstyle{definition}
\newtheorem{definition}[subsection]{Definition}
\newtheorem{example}[subsection]{Example}
\newtheorem{exercise}[subsection]{Exercise}
\newtheorem{situation}[subsection]{Situation}

\theoremstyle{remark}
\newtheorem{remark}[subsection]{Remark}
\newtheorem{remarks}[subsection]{Remarks}

\numberwithin{equation}{subsection}

% Macros
%
\def\lim{\mathop{\mathrm{lim}}\nolimits}
\def\colim{\mathop{\mathrm{colim}}\nolimits}
\def\Spec{\mathop{\mathrm{Spec}}}
\def\Hom{\mathop{\mathrm{Hom}}\nolimits}
\def\Ext{\mathop{\mathrm{Ext}}\nolimits}
\def\SheafHom{\mathop{\mathcal{H}\!\mathit{om}}\nolimits}
\def\SheafExt{\mathop{\mathcal{E}\!\mathit{xt}}\nolimits}
\def\Sch{\mathit{Sch}}
\def\Mor{\mathop{\mathrm{Mor}}\nolimits}
\def\Ob{\mathop{\mathrm{Ob}}\nolimits}
\def\Sh{\mathop{\mathit{Sh}}\nolimits}
\def\NL{\mathop{N\!L}\nolimits}
\def\CH{\mathop{\mathrm{CH}}\nolimits}
\def\proetale{{pro\text{-}\acute{e}tale}}
\def\etale{{\acute{e}tale}}
\def\QCoh{\mathit{QCoh}}
\def\Ker{\mathop{\mathrm{Ker}}}
\def\Im{\mathop{\mathrm{Im}}}
\def\Coker{\mathop{\mathrm{Coker}}}
\def\Coim{\mathop{\mathrm{Coim}}}

% Boxtimes
%
\DeclareMathSymbol{\boxtimes}{\mathbin}{AMSa}{"02}

%
% Macros for moduli stacks/spaces
%
\def\QCohstack{\mathcal{QC}\!\mathit{oh}}
\def\Cohstack{\mathcal{C}\!\mathit{oh}}
\def\Spacesstack{\mathcal{S}\!\mathit{paces}}
\def\Quotfunctor{\mathrm{Quot}}
\def\Hilbfunctor{\mathrm{Hilb}}
\def\Curvesstack{\mathcal{C}\!\mathit{urves}}
\def\Polarizedstack{\mathcal{P}\!\mathit{olarized}}
\def\Complexesstack{\mathcal{C}\!\mathit{omplexes}}
% \Pic is the operator that assigns to X its picard group, usage \Pic(X)
% \Picardstack_{X/B} denotes the Picard stack of X over B
% \Picardfunctor_{X/B} denotes the Picard functor of X over B
\def\Pic{\mathop{\mathrm{Pic}}\nolimits}
\def\Picardstack{\mathcal{P}\!\mathit{ic}}
\def\Picardfunctor{\mathrm{Pic}}
\def\Deformationcategory{\mathcal{D}\!\mathit{ef}}

% --- appended 2026-10-12 (Chief): the project own definitions of the reference/slogan/history
% display environments, copied mechanically from the shared Stacks preamble lines 8-12
% (\usepackage{verbatim} plus the three \newenvironment lines, which the chapter preamble keeps
% verbatim from the project preamble). Nothing added, removed or redefined: these environments
% carry no mathematical content and the project itself typesets them as comments. Without these
% three lines a counted statement whose own text uses them fails with Environment slogan undefined.
\usepackage{verbatim}
\newenvironment{reference}{\comment}{\endcomment}
\newenvironment{slogan}{\comment}{\endcomment}
\newenvironment{history}{\comment}{\endcomment}
\usepackage{cleveref}
\crefname{originallemma}{original lemma}{original lemmas}
\Crefname{originallemma}{Original lemma}{Original lemmas}
\makeatletter
\newlabel{algebra-lemma-directed-colimit-exact}{{00DB}{source}{}{sourcepointer-1}{}}
\newlabel{algebra-lemma-directed-colimit-exact@cref}{{[originallemma][0][]00DB}{[1][0][]source}{}{}{}}
\makeatother
\crefname{originallemma}{original lemma}{original lemmas}
\Crefname{originallemma}{Original lemma}{Original lemmas}
\makeatletter
\newlabel{algebra-lemma-flat-finite-presentation-limit-flat}{{02JO}{source}{}{sourcepointer-2}{}}
\newlabel{algebra-lemma-flat-finite-presentation-limit-flat@cref}{{[originallemma][0][]02JO}{[1][0][]source}{}{}{}}
\makeatother
\crefname{originaltheorem}{original theorem}{original theorems}
\Crefname{originaltheorem}{Original theorem}{Original theorems}
\makeatletter
\newlabel{algebra-theorem-openness-flatness}{{00RC}{source}{}{sourcepointer-3}{}}
\newlabel{algebra-theorem-openness-flatness@cref}{{[originaltheorem][0][]00RC}{[1][0][]source}{}{}{}}
\makeatother
\crefname{originallemma}{original lemma}{original lemmas}
\Crefname{originallemma}{Original lemma}{Original lemmas}
\makeatletter
\newlabel{lemma-quasi-regular-immersion-flat-at-x}{{063T}{source}{}{sourcepointer-4}{}}
\newlabel{lemma-quasi-regular-immersion-flat-at-x@cref}{{[originallemma][0][]063T}{[1][0][]source}{}{}{}}
\makeatother
\crefname{originallemma}{original lemma}{original lemmas}
\Crefname{originallemma}{Original lemma}{Original lemmas}
\makeatletter
\newlabel{lemma-relative-regular-immersion-flat-in-neighbourhood}{{063U}{source}{}{sourcepointer-5}{}}
\newlabel{lemma-relative-regular-immersion-flat-in-neighbourhood@cref}{{[originallemma][0][]063U}{[1][0][]source}{}{}{}}
\makeatother
\crefname{originallemma}{original lemma}{original lemmas}
\Crefname{originallemma}{Original lemma}{Original lemmas}
\makeatletter
\newlabel{more-algebra-lemma-base-change-H1-regular}{{063Q}{source}{}{sourcepointer-6}{}}
\newlabel{more-algebra-lemma-base-change-H1-regular@cref}{{[originallemma][0][]063Q}{[1][0][]source}{}{}{}}
\makeatother
\crefname{originallemma}{original lemma}{original lemmas}
\Crefname{originallemma}{Original lemma}{Original lemmas}
\makeatletter
\newlabel{morphisms-lemma-base-change-module-flat}{{01U8}{source}{}{sourcepointer-7}{}}
\newlabel{morphisms-lemma-base-change-module-flat@cref}{{[originallemma][0][]01U8}{[1][0][]source}{}{}{}}
\makeatother
\title{A relative H1-regular immersion over a locally finitely presented base is flat on a neighbourhood of its image and is a regular immersion stable under base change}
\author{The Stacks Project Authors}
\date{stacks-project:divisors.tex:local-snapshot-20261012; GNU FDL 1.2 or later}
\begin{document}
\maketitle
\noindent\textit{Editorial note.} All delivered mathematical text is the author's own, copied from the pinned Stacks Project chapter file pinned-original-source.tex; no statement or proof step was rewritten, reordered or removed.\par\medskip
\noindent\textit{Editorial note.} Following the multi-claim convention for this source, source\_url is the source-level URL https://stacks.math.columbia.edu/ and the source label lemma-flat-relative-H1-regular is carried in claim\_key and here rather than in a per-claim tag URL.\par\medskip
\noindent\textit{Editorial note.} The preamble is the chapter's shared Stacks Project preamble, retained as sources/stacks-divisors/retained-preamble.tex: a line-for-line excerpt of the shared Stacks preamble staged with this snapshot. No macro was invented or redefined, and this card adds no preamble text of its own.\par\medskip
\noindent\textit{Editorial note.} Statement span 5395-5406 and proof span 5408-5484 come from the halfA pool card operations/pool-cards/stacks-divisors-halfA.json and were re-read in the pinned chapter file before this card was written.\par\medskip
\noindent\textit{Editorial note.} The proof names its approximating rings with the author's own indexing and states the colimit of the Koszul complexes; both are preserved.\par\medskip
\noindent\textit{Editorial note.} The proof contains the sentence 'Assume f\_i is the image of f\_i' for the lifts of the sequence; the author's phrasing is kept unchanged.\par\medskip
\noindent\textit{Editorial note.} 7 cross-references occur inside the counted spans; they are listed in external\_references as navigation pointers only.\par\medskip
\noindent\textit{Allowed standard prerequisites (editorial).}\begin{itemize}
\item The approximation of a flat and finitely presented algebra by a flat algebra of finite type over a finitely generated ring, cited from the algebra chapter. Cited at proof lines 5426-5429. Produces the Noetherian model over which the quotient by the chosen sequence is flat.
\item Surjectivity of the base change map on the first Koszul cohomology of a sequence, and exactness of directed colimits, both cited from the algebra and more-algebra chapters. Cited at proof lines 5437-5439 and 5450-5452. Lets the finitely many generators of the Koszul cohomology of the model be killed after a further finitely generated extension of the base ring.
\item Flatness descends along the base change defining the model, cited from the morphisms chapter. Cited at proof lines 5462-5464. Transfers the flatness statement obtained for the model back to the original morphism.
\item A quasi-regular immersion is flat at each point of its image, and the flat locus of a finite type morphism of Noetherian schemes is open, cited from the same chapter and the algebra chapter. Cited at proof lines 5468-5473. Supplies pointwise flatness in the reduced situation and turns it into an open neighbourhood.
\item A relative regular immersion which is flat in a neighbourhood of its image is a regular immersion with the property stable under base change, cited from the same chapter. Cited at proof lines 5474-5476. Gives part (2) once part (1) is known.
\end{itemize}\medskip
\ifdefined\hypertarget\hypertarget{sourcepointer-1}{}\fi
\noindent\textit{Source reference.} Original lemma 00DB: focused-seeds/AGP/sources/stacks-snapshot-20261011/tags.txt, source lines 491--491. Same-chapter reference used in the retained author text at source line(s) 5458. It is supplied only as a navigation pointer so that the retained text compiles; the counted item delivers the author's own statement and the author's complete proof, and no part of the referenced item is copied in or substituted for the counted argument.\par\medskip
\ifdefined\hypertarget\hypertarget{sourcepointer-2}{}\fi
\noindent\textit{Source reference.} Original lemma 02JO: focused-seeds/AGP/sources/stacks-snapshot-20261011/tags.txt, source lines 3343--3343. Same-chapter reference used in the retained author text at source line(s) 4106, 4234, 5430. It is supplied only as a navigation pointer so that the retained text compiles; the counted item delivers the author's own statement and the author's complete proof, and no part of the referenced item is copied in or substituted for the counted argument.\par\medskip
\ifdefined\hypertarget\hypertarget{sourcepointer-3}{}\fi
\noindent\textit{Source reference.} Original theorem 00RC: focused-seeds/AGP/sources/stacks-snapshot-20261011/tags.txt, source lines 998--998. Same-chapter reference used in the retained author text at source line(s) 4132, 4252, 5481. It is supplied only as a navigation pointer so that the retained text compiles; the counted item delivers the author's own statement and the author's complete proof, and no part of the referenced item is copied in or substituted for the counted argument.\par\medskip
\ifdefined\hypertarget\hypertarget{sourcepointer-4}{}\fi
\noindent\textit{Source reference.} Original lemma 063T: focused-seeds/AGP/sources/stacks-snapshot-20261011/tags.txt, source lines 7842--7842. Same-chapter reference used in the retained author text at source line(s) 5479. It is supplied only as a navigation pointer so that the retained text compiles; the counted item delivers the author's own statement and the author's complete proof, and no part of the referenced item is copied in or substituted for the counted argument.\par\medskip
\ifdefined\hypertarget\hypertarget{sourcepointer-5}{}\fi
\noindent\textit{Source reference.} Original lemma 063U: focused-seeds/AGP/sources/stacks-snapshot-20261011/tags.txt, source lines 7843--7843. Same-chapter reference used in the retained author text at source line(s) 5301, 5387, 5483, 5511. It is supplied only as a navigation pointer so that the retained text compiles; the counted item delivers the author's own statement and the author's complete proof, and no part of the referenced item is copied in or substituted for the counted argument.\par\medskip
\ifdefined\hypertarget\hypertarget{sourcepointer-6}{}\fi
\noindent\textit{Source reference.} Original lemma 063Q: focused-seeds/AGP/sources/stacks-snapshot-20261011/tags.txt, source lines 7839--7839. Same-chapter reference used in the retained author text at source line(s) 4113, 5441. It is supplied only as a navigation pointer so that the retained text compiles; the counted item delivers the author's own statement and the author's complete proof, and no part of the referenced item is copied in or substituted for the counted argument.\par\medskip
\ifdefined\hypertarget\hypertarget{sourcepointer-7}{}\fi
\noindent\textit{Source reference.} Original lemma 01U8: focused-seeds/AGP/sources/stacks-snapshot-20261011/tags.txt, source lines 2401--2401. Same-chapter reference used in the retained author text at source line(s) 4126, 5474. It is supplied only as a navigation pointer so that the retained text compiles; the counted item delivers the author's own statement and the author's complete proof, and no part of the referenced item is copied in or substituted for the counted argument.\par\medskip
\begin{lemma}
\label{lemma-flat-relative-H1-regular}
Let $X \to S$ be a morphism of schemes.
Let $Z \to X$ be a relative $H_1$-regular immersion.
Assume $X \to S$ is locally of finite presentation. Then
\begin{enumerate}
\item there exists an open subscheme $U \subset X$ such that
$Z \subset U$ and such that $U \to S$ is flat, and
\item $Z \to X$ is a regular immersion and the same remains
true after any base change.
\end{enumerate}
\end{lemma}

\begin{proof}
Pick $x \in Z$. To prove (1) suffices to find an open neighbourhood
$U \subset X$ of $x$ such that $U \to S$ is flat. Hence the lemma reduces
to the case that $X = \Spec(B)$ and $S = \Spec(A)$ are affine
and that $Z$ is given by an $H_1$-regular sequence $f_1, \ldots, f_r \in B$.
By assumption $B$ is a finitely presented $A$-algebra and
$B/(f_1, \ldots, f_r)B$ is a flat $A$-algebra. We are going to use
absolute Noetherian approximation.

\medskip\noindent
Write $B = A[x_1, \ldots, x_n]/(g_1, \ldots, g_m)$. Assume
$f_i$ is the image of $f_i' \in A[x_1, \ldots, x_n]$. Choose a finite type
$\mathbf{Z}$-subalgebra $A_0 \subset A$ such that all the coefficients
of the polynomials $f_1', \ldots, f_r', g_1, \ldots, g_m$ are in $A_0$.
We set $B_0 = A_0[x_1, \ldots, x_n]/(g_1, \ldots, g_m)$ and we denote
$f_{i, 0}$ the image of $f_i'$ in $B_0$. Then $B = B_0 \otimes_{A_0} A$
and
$$
B/(f_1, \ldots, f_r) =
B_0/(f_{0, 1}, \ldots, f_{0, r}) \otimes_{A_0} A.
$$
By
Algebra, Lemma \ref{algebra-lemma-flat-finite-presentation-limit-flat}
we may, after enlarging $A_0$, assume that
$B_0/(f_{0, 1}, \ldots, f_{0, r})$ is flat over $A_0$.
It may not be the case at this point that the Koszul cohomology group
$H_1(K_\bullet(B_0, f_{0, 1}, \ldots, f_{0, r}))$ is zero.
On the other hand, as $B_0$ is Noetherian, it is a finitely
generated $B_0$-module. Let
$\xi_1, \ldots, \xi_n \in H_1(K_\bullet(B_0, f_{0, 1}, \ldots, f_{0, r}))$
be generators. Let $A_0 \subset A_1 \subset A$ be a larger finite type
$\mathbf{Z}$-subalgebra of $A$. Denote $f_{1, i}$ the image
of $f_{0, i}$ in $B_1 = B_0 \otimes_{A_0} A_1$. By
More on Algebra, Lemma \ref{more-algebra-lemma-base-change-H1-regular}
the map
$$
H_1(K_\bullet(B_0, f_{0, 1}, \ldots, f_{0, r})) \otimes_{A_0} A_1
\longrightarrow
H_1(K_\bullet(B_1, f_{1, 1}, \ldots, f_{1, r}))
$$
is surjective. Furthermore, it is clear that the colimit (over all
choices of $A_1$ as above) of the
complexes $K_\bullet(B_1, f_{1, 1}, \ldots, f_{1, r})$ is the complex
$K_\bullet(B, f_1, \ldots, f_r)$ which is acyclic in degree $1$. Hence
$$
\colim_{A_0 \subset A_1 \subset A}
H_1(K_\bullet(B_1, f_{1, 1}, \ldots, f_{1, r}))
= 0
$$
by
Algebra, Lemma \ref{algebra-lemma-directed-colimit-exact}.
Thus we can find a choice of $A_1$ such that $\xi_1, \ldots, \xi_n$
all map to zero in $H_1(K_\bullet(B_1, f_{1, 1}, \ldots, f_{1, r}))$.
In other words, the Koszul cohomology group
$H_1(K_\bullet(B_1, f_{1, 1}, \ldots, f_{1, r}))$
is zero.

\medskip\noindent
Consider the morphism of affine schemes
$X_1 \to S_1$ equal to $\Spec$ of the
ring map $A_1 \to B_1$ and
$Z_1 = \Spec(B_1/(f_{1, 1}, \ldots, f_{1, r}))$.
Since $B = B_1 \otimes_{A_1} A$, i.e., $X = X_1 \times_{S_1} S$,
and similarly $Z = Z_1 \times_S S_1$,
it now suffices to prove (1) for $X_1 \to S_1$ and the relative
$H_1$-regular immersion $Z_1 \to X_1$, see
Morphisms, Lemma \ref{morphisms-lemma-base-change-module-flat}.
Hence we have reduced to the case where $X \to S$ is a finite type
morphism of Noetherian schemes.
In this case we know that $X \to S$ is flat at every
point of $Z$ by
Lemma \ref{lemma-quasi-regular-immersion-flat-at-x}.
Combined with the fact that the flat locus is open in this case, see
Algebra, Theorem \ref{algebra-theorem-openness-flatness}
we see that (1) holds. Part (2) then follows from an application of
Lemma \ref{lemma-relative-regular-immersion-flat-in-neighbourhood}.
\end{proof}

\end{document}
